\section*{Chapter \ref{ch:m1:index-construction-and-rebalancing} --- Index Construction and Rebalancing}

\begin{solution}{exo:m1:index-construction-and-rebalancing:1}
Total capitalisation \$20 billion; float-adjusted $20 \times 0.65 = \$13$
billion; weight $13/5\,000 = 0.26\%$.
\end{solution}

\begin{solution}{exo:m1:index-construction-and-rebalancing:2}
$\varphi = 2/5 = 40\%$ of every member's float. Demand: $0.40 \times 13 =
\$5.2$ billion, thirteen days of volume.
\end{solution}

\begin{solution}{exo:m1:index-construction-and-rebalancing:3}
Member at 108: stays ($108 \le 110$). Outsider at 95: does not enter (it
needs 90 or better), unless places remain to be filled by rank. Outsider at
88: enters. Member at 112: deleted.
\end{solution}

\begin{solution}{exo:m1:index-construction-and-rebalancing:4}
In the large index: membership was fixed on 30 April, and later prices do not
change it. The fall changes the \emph{size} of everything: the stock enters
with a weight and a tracker demand 30\% smaller in dollars (the same number
of shares, since trackers hold a fraction $\varphi$ of the float). The firm
that bought in April was right about the event and lost 30\% on the stock: an
index position is a position in the stock first. The event's few percent
cannot pay for unhedged exposure to the company.
\end{solution}

\begin{solution}{exo:m1:index-construction-and-rebalancing:5}
$10\,000 \times 1.5\% = \$150$ billion. The 6\% member becomes $6\% \times
0.985 = 5.91\%$: a sale of $0.09\% \times 10\,000 = \$9$ billion. Two steps
would halve the trade at each close and the strain on liquidity; but each
step is a separate event to be anticipated, the index would hold a
half-weight that represents nothing for a week, and trackers would bear
tracking decisions twice. One step on the most liquid close available
concentrates the liquidity where the demand is.
\end{solution}

\begin{solution}{exo:m1:index-construction-and-rebalancing:6}
One day: $0.7 \times 2.5\% \times \sqrt{13} = 6.3\%$, the order of the 7.4\%
measured in the 1990s (and the rule is being used far outside its range).
Spread over twenty days: $0.7 \times 2.5\% \times \sqrt{13/20} = 1.4\%$ a day
of temporary impact, by participants who then hold the shares at risk for
weeks; competition between them, and the supply from the index the stock
leaves, bring the visible effect down towards the 0.3\% of the last decade.
The same demand, absorbed by a longer and more crowded queue.
\end{solution}

\begin{solution}{exo:m1:index-construction-and-rebalancing:7}
$b = 0$: 18.5 additions per review, one-way turnover 1.92\%. $b = 30$: 9.2
additions, 1.18\%. The fund trades twice the one-way turnover; the saving is
$2 \times 0.74\% \times 15$ basis points $= 0.22$ basis point a year.
Explicit costs are not the reason for buffers: the names that flip at the
boundary are the smallest in the index. What buffers reduce is the number of
predictable events in which trackers pay impact that the 15 basis points do
not capture.
\end{solution}

\begin{solution}{exo:m1:index-construction-and-rebalancing:8}
(i) The effect measured on that decade has since fallen to a few tenths of a
percent: the backtest period is the hypothesis. (ii) The 5\% was measured
from prices at which a small study could trade; every participant now buys
at the announcement, so the next open already contains the move, and
twenty-five events do not have unlimited capacity. (iii) Returns per event
do not add to an annual return: events overlap, capital is tied for days to
weeks, the positions carry unhedged stock risk far larger than the edge, and
costs were ignored. One may add that many additions are now predicted
before they are announced, which moves the return outside the window
traded.
\end{solution}

\begin{solution}{pb:m1:index-construction-and-rebalancing:1}
\textbf{1.} $200/3\,000 = 6.67\%$; $4\,000/50\,000 = 8\%$.
\textbf{2.} $6.67\% \times 1.2$ billion $= \$80$ million; 3.33 million shares.
\textbf{3.} 3.7 days.
\textbf{4.} Sale $6.67\% \times 9 = \$600$ million; purchase $8\% \times 9 =
\$720$ million; net $+\$120$ million.
\textbf{5.} The sale is 4.0 days of volume, the net 0.8 day. If both groups
trade in the same closing auction only the net presses on the price: a
migration is mostly a cross between two sets of trackers.
\textbf{6.} 2 million shares against a normal auction of 72\,000: 28 times.
\textbf{7.} $0.7 \times 3\% \times \sqrt{3.7} = 4.0\%$.
\textbf{8.} $0.7 \times 3\% \times \sqrt{3.7/20} = 0.9\%$ a day.
\textbf{9.} The anticipators, who bought over the preceding weeks from
ordinary sellers at a rate the market could bear, and \omterm{def:m1:what-a-trading-firm-does:market-maker}{market makers} who go
short into the auction and buy back in the following days. The auction's
imbalance feed (\cref{ch:m1:auctions}) tells them how much is still missing.
\textbf{10.} $0.8 \times 3\% - 0.2 \times 2\% - 0.4\% = 1.6\%$: \$80\,000.
\textbf{11.} The return is 3\% or $-2\%$: standard deviation $\sqrt{0.8
\times 0.2} \times 5\% = 2\%$, \$100\,000, before any ordinary stock risk.
\textbf{12.} \$2 million, \$500\,000, ratio 4.0.
\textbf{13.} Standard deviation $100\,000 \times \sqrt{25 + 600 \times 0.3} =
\$1.43$ million; ratio 1.4.
\textbf{14.} $p = (0.4\% + 2\%)/5\% = 0.48$.
\textbf{15.} All the positions are long small stocks near one breakpoint: a
market move before the rank day moves the breakpoint and all the
probabilities together; and the same anticipators hold all the names, so
when one of them reduces risk every position falls at once.
\textbf{16.} Each event moves fewer names and smaller weights, since half as
much drift accumulates between reviews; there are twice as many dates. The
total to be traded falls somewhat (fewer stocks wander far from their index),
and the business becomes less seasonal and less concentrated on one Friday.
\textbf{17.} They are judged on \omterm{def:m1:the-buy-side:benchmark}{tracking error} against an index computed at
that close. Buying early at a better price is a gain of a few basis points
if it works and an unexplainable deviation if it does not.
\textbf{18.} The fund's investors, through \omterm{def:m1:the-buy-side:benchmark}{tracking error} in both
directions; they find out from the tracking difference in the annual report,
and from nothing else.
\textbf{19.} \textbf{3.7 days of average volume.}
\textbf{20.} They are the warehouse that turns a one-minute demand into a
three-week one, and their profit is the fee for that.
\end{solution}

\begin{solution}{iq:m1:index-construction-and-rebalancing:1}
Trackers must buy a known fraction of its float at the effective close.
Historically the stock jumped at the announcement and rose into the effective
date, with a partial reversal. For the large US index the measured
announcement effect has gone from several percent in the 1990s to almost
nothing: the move now happens earlier, as additions are predicted, and the
demand is met by firms that position in advance and by the trackers of the
index the stock leaves.
\iqlookfor{the timeline (prediction, announcement, effective close), and
knowledge that the textbook effect has shrunk.}
\end{solution}

\begin{solution}{iq:m1:index-construction-and-rebalancing:2}
An index is meant to be holdable by everyone simultaneously. With full
capitalisation weights, trackers would need shares that a founder or a state
will never sell: for a company 80\% held by its founder, a demand five times
larger relative to the tradable shares than for its neighbours, a squeeze at
every inflow, and a price distorted by the index itself.
\iqlookfor{the ``everyone at once'' argument and a concrete squeeze.}
\end{solution}

\begin{solution}{iq:m1:index-construction-and-rebalancing:3}
$8\% \times 10$ billion $= \$800$ million. Whether it is a lot depends on the
stock's daily volume (one day or fifteen), on the size of its closing
auction, on whether it leaves another index with trackers of its own (the
net, not the gross, matters), and on how much has already been bought by
anticipators, which the run-up and the short interest of the auction will
show.
\iqlookfor{the fraction-of-float shortcut and the net of migration.}
\end{solution}

\begin{solution}{iq:m1:index-construction-and-rebalancing:4}
Rebuild the provider's universe and its total capitalisations from the
methodology; rank; locate the breakpoints and apply the bands to current
members; before the rank day attach a crossing probability to each boundary
name from distance and volatility. Errors: share counts (multiple classes,
recent issues, filings the provider reads differently), eligibility details
(domicile, minimum float, listing), the breakpoints themselves moving with
the market, and corporate actions between rank day and effective date.
\iqlookfor{share-count humility and the moving breakpoint.}
\end{solution}

\begin{solution}{iq:m1:index-construction-and-rebalancing:5}
The measured effect is the return in a window that starts at the
announcement. More passive assets mean a larger demand, but also a larger
reward for predicting and supplying it: the price adjusts before the window,
more capital competes to be the seller at the close, and many additions
migrate from a sister index whose trackers are natural sellers. The shock
grew; its price fell.
\iqlookfor{the distinction between the demand and the measured return.}
\end{solution}

\begin{solution}{iq:m1:index-construction-and-rebalancing:6}
Risks: prediction error (the name is not added); stock risk between entry and
the close, which dwarfs the edge and must be hedged by sector and size;
crowding, visible as correlated drawdowns of all candidates when a
competitor de-risks; breakpoints moving with the market; rule changes by the
provider; and execution in one auction. Sizing: by days of volume to exit,
not by dollars; by the probability-weighted edge net of hedging cost; with a
cap on the share of the expected tracker demand that the book may hold, since
beyond some share the book is its own exit.
\iqlookfor{hedging of the stock leg, crowding, and sizing in units of
liquidity.}
\end{solution}
